% methods_def.tex -- Section III-C (sensitivity), III-D (corrected baseline), III-E (wearable).
% Input on the last line of methods.tex. Labels: sec:sens, sec:corrected, sec:wearable.
% Sources: S1 grid and limit (A1-1, A1-3); 1.3 cm pelvis lift, joints fixed (M36, PS6, PS7);
%   foot points (M29); S3 deployment form 20 frames at 100 Hz (A3-34); reference motion 6.45 s at
%   100 Hz (GT6); its SO+JR labels are unusable (label_mask keeps 0 right-knee frames; sep23/README).
%   Corrected baseline = user decision (a)/(b), 2026-11-10. Training set = train_mlp *_sojrqc input
%   sets: keep_r of <trial>/jcf_output_sojr/label_mask.npz (fixed right-foot COP trials, then the six
%   SO/dynamics frame rules; the 5 BW flag is stored, NOT applied), carried to the 100 Hz grid with
%   LABEL_MASK_TOL = 0.05, physical M_x sign (MX_SIGN_PHYSICAL). 25.4 % = kept_stance_frac_r of the
%   training walking pool in results/thesis/labels_qc/label_masks/by_lab.csv (adding the 5 BW flag
%   would remove another 2.5 % of the kept stance time, fy5_after_so_frac_of_kept_stance_r). The
%   legacy trial filters (6 BW, 0.15 BW*H) would drop 86 % of the kept training frames
%   (label_masks/dry_build_fpopt/trial_filters.csv), so the corrected runs switch them off
%   (--max-peak 1e9 --max-moment 1e9); keep in sync with the training commands.
%   Lencioni = the lab that earlier drafts called 'Tiziana' (data folder name). S2 pelvis lift: the numbers are [4]'s own, cited as such in results.tex; we did not repeat the
%   test (no checkpoint was lift-tested). S2 pelvis-angle conversion and S4: results/thesis/sep23
%   (converted pelvis angles are the primary inputs). S3: legacy in vivo, audit_gc_legacy_native vs
%   audit_gc. Seed 2 best epoch 8 (1-based), val 0.363; seeds 1/3 0.300/0.305 (sep23/README.md).
%   Simulated suit: two held-out trials, noise 0/1/2 deg and 1 deg + 0.5 deg/s drift (pn3_sim/README).
%   Wearable: A3-01..A3-03, A3-05, A3-22..A3-25, A3-29..A3-34, A3-39, A3-41, A3-07, A3-08;
%   Wang et al. data: results/thesis/twente/EXPLORE_NOTES.md sections 1-6 (9 people, 54 treadmill
%   trials at 0.25-1.50 m/s, 8 sensors, optical IK, treadmill force, separate static pose, no
%   knee-load labels).
\subsection{Sensitivity Protocol}\label{sec:sens}
\textit{S1, constants.} We recompute $\Fmed$ for $\phi \in \{0.4, 0.5, 0.6\}$ and
$d \in \{40, 45, 50\}$~mm, with the limit at 1.4~BW, and report the mean absolute error (MAE) and
the recall and precision of limit crossings. \textit{S2, inputs.} We swap the foot points and, on
the reference motion, compare two pelvis-angle conventions. \textit{S3, alignment.} Legacy models
are scored in vivo on training-style and on aligned 100~Hz inputs, as the optimizer and IMU
driver supply. \textit{S4, training run.} Models differing only in initialization seed are
scored on the test split, in vivo, and on the optimizer's reference motion~\citeChong{},
which lacks usable labels.

\subsection{Corrected Baseline}\label{sec:corrected}
The corrected baseline keeps the legacy architecture and recipe but (i)~uses SO+JR labels,
(ii)~uses the physical sign in the $\Fmed$ loss and stores it in the checkpoint, (iii)~builds
inputs and labels on the JR time grid resampled to 100~Hz, so a window spans 200~ms as in
deployment, and (iv)~trains only on frames that pass the quality rules (Table~\ref{tab:qc}),
which replace the legacy trial filters; a 100~Hz frame passes when at most 5\% of its
interpolated label comes from flagged frames. Seeds 1--3 use the optimizer's foot point and one
run (seed~1) the calcaneus origin, all on one person split. Seed~2 reached its best validation
loss at epoch~8 (0.363, against 0.300 and 0.305); we keep it. Legacy checkpoints keep their
formula, so results computed with them stay reproducible.

\subsection{Wearable Path}\label{sec:wearable}
A Perception Neuron~3 suit streams fused orientations of 8 sensors (pelvis, torso, thighs,
shanks, feet) at 60~Hz, but no pelvis position. After a 1-s standing calibration, an unscented
Kalman filter (UKF)~\cite{julier2004unscented} with $2\times38$ cubature
points~\cite{arasaratnam2009cubature} estimates 19 coordinates of a height-scaled Rajagopal model
and their velocities from orientation residuals ($\sigma = 2^\circ$), and legged
odometry~\cite{bloesch2012state} holds each planted foot ($\sigma = 1.5$~cm). Hysteresis detectors
on foot angular speed~\cite{skog2010zero} set the contact flags; as they switch on at foot flat,
about 90~ms after heel strike, each onset is back-dated by up to 0.2~s to the speed minimum at
initial contact. A 20-frame, 100~Hz window feeds the surrogate (90~ms lag); results go over Redis
to a visualizer.

We test the path on a simulated suit (orientations from two held-out trials with 0, 1 and
2$^\circ$ noise, and 1$^\circ$ with 0.5$^\circ$/s heading drift) and on the Xsens recordings of
Wang et al.~\cite{wang2023twente}: 9 people, 54 treadmill trials at 0.25--1.50~m/s, 8 sensors at
the suit's locations, optical motion capture, treadmill force and a separate static calibration
pose. Lacking knee-load labels, we score the IMU path against optical inverse kinematics (IK) and
against the same surrogate run on optical kinematics.
